\documentclass[11pt]{article}
\usepackage{amsmath,amssymb,amsthm,fullpage}
\newtheorem{theorem}{Theorem}
\title{Star Covers for Unbalanced Exceptional Faces}
\author{}
\date{November 23, 2026}
\begin{document}
\maketitle

\paragraph{Status.}
The face approximation and concrete compiler transfer below are Lean
checked in \texttt{AxisZeroStarFaces.lean} and
\texttt{AxisZeroUnbalancedTransfer.lean}. They use the unchanged original
atoms and invariants. The unbalanced asymptotic calculation is a complete
written implication, but its local compiler premise is unproved. This
note does not establish either accepted asymptotic target.

Write $K(m)$ for the existing \texttt{monomialStarCoverNumber m}. Its
checked bound is $K(m)\leq 11\,2^m/m$ for $m\geq1$.
An upward monomial star with center $U\subseteq[m]$ consists of $U$ and
some of its one-element extensions. Thus $K(m)$ stars cover all subsets
of $[m]$.

\begin{theorem}[Real-valued face correction]
Fix disjoint coordinate sets $S$ and $J$ in the original $n$-cube, with
$|J|=m$. Every function
\[
 x\longmapsto x_S g(x_J),\qquad x_S=\prod_{i\in S}x_i,
 \qquad g:\{0,1\}^J\longrightarrow\mathbb R,
\]
is uniformly approximable, to arbitrary positive accuracy on the whole
cube, by a bias plus $K(m)$ original admissible atoms.
\end{theorem}
\begin{proof}
Expand $g$ in squarefree monomials,
$g(x_J)=\sum_{U\subseteq J}c_Ux_U$. Assign every $U$ to a covering
upward star with center $V$. The terms assigned to $V$ sum to
$x_V L_V(x_J)$ for an affine $L_V$. Multiplying by $x_S$ gives
$x_{S\cup V}L_V(x_J)$, since $S$ and $J$ are disjoint. The existing
one-atom approximation theorem for a monomial times an affine function
applies to each such term. Approximate each of the $K(m)$ terms closely
enough that their total error is less than the prescribed accuracy.
Constants are included in this grouping; no subtraction of a constant
outside the face is used.
\end{proof}

\begin{theorem}[Two faces]
Suppose a score $G$ with $H$ original heads has the correct nonzero
strict signs outside the union of two positive coordinate faces of
free dimensions $a,b$. Then the whole Boolean function has an original
representation with at most $H+K(a)+K(b)$ heads.
\end{theorem}
\begin{proof}
Let $T(x)$ equal $1$ or $-1$ according to the desired Boolean value.
For the first face, let $P$ be its fixed-coordinate monomial and let
$\rho(x)$ set its fixed coordinates to one. The correction
$R(x)=P(x)(T(\rho(x))-G(\rho(x)))$ replaces $G$ by $T$ on that face
and vanishes elsewhere. Write $G_1=G+R$. For the second face use
$S(x)=Q(x)(T(\tau(x))-G_1(\tau(x)))$.
Then $G_1+S$ has correct strict signs everywhere. On the intersection,
$G_1=T$, so the second correction is zero. The first theorem approximates
$R$ and $S$ with $K(a)$ and $K(b)$ heads. The finite cube has a positive
minimum strict margin for $G_1+S$. Choose the total approximation error
smaller than this margin. This also preserves every previously obtained
sign away from the faces.
\end{proof}

\begin{theorem}[Unbalanced concrete transfer]
Let $a,b\geq1$. Suppose every $a$ by $b$ sign table admits a reverse
axis-zero local compiler with at most $M$ heads, positive denominators
on the full $(a+b)$-cube, and strictly negative slopes, as required by
\texttt{HasReverseAxisCompiler a b M}. Then
\[
 W(a+b)\leq 484\frac{2^{a+b}M}{ab}
       +K(a)+K(b)
 \leq484\frac{2^{a+b}M}{ab}
       +11\frac{2^a}{a}+11\frac{2^b}{b}.
\]
\end{theorem}
\begin{proof}
Use the checked open upper-star covers of the two coordinate blocks,
with at most $22\,2^a/a$ and $22\,2^b/b$ centers. The checked descending
cell enumeration and reverse gluing theorem use one $M$-head module per
pair of centers, and give strict correct signs outside the two positive
exceptional faces. Apply the preceding theorem. All biases are absorbed
by the existing gluing construction before this correction; this argument
introduces no extra head for them. The exact integer statement is
$H^*(f)\leq |D\times E|M+K(a)+K(b)$. Maximizing over the unchanged
Boolean-function domain yields the displayed bound.
\end{proof}

\paragraph{Why the unbalanced split is now viable.}
Put $b=\lfloor\log_2 n\rfloor$ and $a=n-b$. For all sufficiently large
$n$, we have $a\geq n/2$ and $n/2<2^b\leq n$. Consequently
\[
 \frac{2^a}{a}\leq\frac{4\,2^n}{n^2},\qquad
 \frac{2^b}{b}\leq\frac{n}{b}.
\]
If the missing compiler could be proved with
$M(a,b)\leq Cb/\log(b+1)$ along these splits, the main term would be
at most $968C\,2^n/(n\log(b+1))$, while both face terms are smaller
asymptotically. Since $\log(b+1)\sim\log\log n$, this would prove
target (A), with natural logarithms. This is a strictly identified
unproved hypothesis, not a claimed compiler theorem. In particular,
a VC lower bound without the compact full-cube denominator condition
would not by itself discharge the hypothesis.

\paragraph{Checked declarations.}
In namespace \texttt{HeadComplexity.AxisZeroGluing}, the central new
results are
\texttt{positive\_face\_uniformlyAtomsApproximable\_star},
\texttt{repair\_two\_positive\_faces\_star},
\texttt{HStar\_le\_product\_cover\_of\_compiler\_star\_faces}, and
\texttt{worstCaseHeadComplexity\_real\_le\_of\_reverse\_axis\_compiler\_star\_faces\_eleven}.
The original definitions and earlier transfer results are unchanged.
\end{document}
